%There is value is teamwork. Having competed the challenge in both individually and in different  teams over the years, there are pros and cons to both. A good team can lift your ability, not so great teammates just leave it when challenges get tough, but that’s fine.. at leats you know you’ want stick with it. I’ve experienced both sides. Having a supportive team as well as one that just left and didn’t do anything. 
%Also wanna talk about how I didn’t do any programming, but being part of a team meant I could learn from others. It’s a good team building exercise, good way to know people, especially if you’ve ended up changing schools for sixth form or something which is what I ended up doing. Had to form a new team. Something about the experience of shared suffering… people bringing different set of skills which can be useful. Some are very good ate rogramming.. others at recognising what works and doesn’t… good to bounce back ideas form each other.. but also can lead to argument and falling out… that’s tough. So just, look after each other. 
%Other resources you could have as a team…we had a shared GitHub reporsitoeyr we worked on over the summer and writing solutions to ciphers that were common in the competition. 
%Team captain login- remember that only that account that submit solutions, so if you team captain wants it might be a good idea to share that with the team so that others can submit the solution in case they can’t get hold of the captain. 


\section{Methods of Competing}

Having competed in the competition both individually and as part of different teams overs, I've had some experience of both sides. 

Working alone can be focused as you move at your own pace, follow your own ideas and don't have to coordinate with anyone else. But it can also be isolating at times, especially when you get stuck. 

A good team on the other hand can lift your ability. It gives your different perspectives, methods of thinking, and people to bounce ideas off. At the same time, not every works efficiently and sometimes when challenges get tough, people drift away. So if you ever find yourself left alone towards the end when you started of with a team, rest assured, it's quite normal. From my experience, being the only left taught me the kind of person I wanted to be in a team, namely someone who stays, contributes, and stick with it when things get tough.

Having joined a new school for sixth form, I had the fun change of joining a new team from scratch. Working together on the cipher challenge later turned out to be one of the easiest ways to get to know people. It's genuinely an amazing experience to work together to solve a problem that no one fully understands. 

\section{Learning from Each Other}

One of the biggest advantages of working in a team is the opportunity to learn from one another. 

For instance, when I first started competing in the NCC, I could barely code. But working alongside friends who shared this enthusiasm for problem-solving created an environment in which I saw my confidence in coding improve. At the start of my sixth form year, I had just over a year of coding experience and could identify most commonly used ciphers in the challenge, but I didn't really see myself as the `coding person'\footnote{The only Python programs I recall writing were an IOC calculator and a simple Vigenère cipher solver for a known key}. Yet, through the process of contributing to our team's shared GitHub repository, I found myself working on annealing and modified Playfair ciphers. Whether it was optimising the hill climb parameters for various programs, or spotting a silly inequality error in the quadgram evaluation block, the shared teamwork made even the most daunting ciphers feel achievable.

At the same time, programming isn't the only skill that matters. Some people are amazing at coding and automating processes, whereas others are better at spotting patterns, recognising when something doesn't quite it, and knowing when to move onto another idea. 

A good team brings all of these strengths together.

\section{Shared Thinking}

An idea that seems convincing in your own head can start to fall apart when you try to explain it out loud. On the other hand, someone else might take the same idea and see something you didn't. There's a lot of scope for progress from conversations with teammates. 

There's also something to be said for the experience itself. The shared frustration, the moments where nothing makes sense, and the occasional breakthrough when something suddenly clicks. That shared experience is a big part of what makes this competition memorable. 

Teamwork is not always smooth. People work and think at different speeds and sometimes ideas clash and disagreements happen. When things get difficult, it's easy for communication to break down. 

It helps to remember that everyone is working towards the same common goal. So taking a step back, listening to each other and being willing to let go of your own initial ideas to intake new information can make a big difference. 

\section{Working Practically}

Listed below are some things that helped our team work more efficiently, and I hope you might find something useful here: 

\begin{itemize}
    \item Sharing Progress - it can be useful to have something like a team Discord server or any group chat so that you can communicate and keep each other updated with progress 
    \item Divide and Explore - different people can test different ideas at the same time which can speed up the process of ruling things out 
    \item Shared resources - my team kept a shared GitHub repo where we stored all our scripts and solutions to most common ciphers. Having  a shared codebase can be very useful in the long term and make later challenges easier to approach  
    \item Team captain login - since only the 'Team Captain' account can submit solutions, it might be worth to make sure more than one person has access to it in case the team captain is unavailable 
\end{itemize}
